% Physical unit 028; accepted Session 03 after fresh 400-dpi authority replay.
\GordanUnitStart{028}{printed page 23}
\GordanTopFolio{23}
\begingroup\small
\caput{VI}
\sect{1}

We have seen that the most important moments of the geodetic line can be deduced from the formulae of Chapter II proposed by the illustrious Legendre.

For computing lengths and extensions from $l$, however, the formulae developed by the illustrious Jacobi in the booklet
\[
\text{\gq{Solution nouvelle d'un problème de Géodésie}}
\]
are far more suitable, by passing from elliptic integrals to elliptic functions.

\sect{2}

Following him, in place of the angle $\varphi$, which is variable by itself, we introduce the variable
\[
 u=\int_0^\varphi {d\varphi\over \sqrt{1-x^2\sin^2\varphi}}
\]
and express all quantities by means of its functions. Accordingly, for computing $\varphi$ we use the formulae arising from the inverse equation of the integral; but for computing the remaining quantities it is necessary to carry these values back to the equations of Chapter II.

Thus there results:
\[
\begin{gathered}
 \varphi=\operatorname{am}u,\qquad \sin\varphi=Su,
 \qquad \cos\varphi=Cu,\qquad \sqrt{1-x^2\sin^2\varphi}=\Delta u,\\
 \sin l=\sin l_0\,Cu,\qquad
 \cos l=\cos l_0\sqrt{1+\tg^2 l_0\,S^2u},\\
 \tg l=\tg l_0{Cu\over \sqrt{1+\tg^2l_0\,S^2u}},\\
 \sqrt{1-e^2\cos^2l}={b\over a}\cdot {\Delta u\over x'}
 =\sqrt{1-e^2\cos^2l_0}\,\Delta(u),\\
 \tg B=\tg B_0{Cu\over \sqrt{1+\tg^2l_0\,S^2u}},\\
 \sin B=\sin B_0{Cu\over \Delta u},
 \qquad
 \cos B=\cos B_0{\sqrt{1+\tg^2l_0\,S^2u}\over \Delta u},\\
 \sqrt{1-e^2\sin^2B}={\cos B\over \cos l}={x'\over \Delta u}.
\end{gathered}
\]
\endgroup
